% !TEX root = ../../main.tex

\subsection{Sibling Test: Are the Two Children Distinct?}

The sibling test is the second inferential step. For each parent $u$ with
exactly two children $l(u)$ and $r(u)$, it asks whether the two children are
genuinely different from one another rather than merely both differing from the
parent. The edge test alone is not enough to justify keeping a split. A parent
can contain one deviating branch, or a weak left-right asymmetry, without
supporting two stable child clusters. The sibling test supplies that missing
question: once a branch has left its ancestral context, are the two children
actually distinct enough to become separate clusters?
Because the hierarchy is a rooted binary tree, this sibling relation has a
fixed form: every non-root node has exactly one parent and therefore at most one
sibling, namely the other child of that same parent. The sibling test is thus
always a single pairwise comparison, not a comparison against a larger sibling
set.

\paragraph{Sibling null and pooled baseline.}
Formally, it tests the sibling null
\begin{equation}
H_0^{\mathrm{sib}}(u): \theta_{l(u)} = \theta_{r(u)}.
\end{equation}
For Bernoulli coordinates, the pooled null proportion is
\begin{equation}
\bar{\theta}_{u,j}
=
\frac{n_{l(u)}\theta_{l(u),j} + n_{r(u)}\theta_{r(u),j}}
{n_{l(u)} + n_{r(u)}}.
\end{equation}

\paragraph{Standardized sibling contrast.}
Because the two children are disjoint subtrees under the same parent, the
coordinatewise variance takes the usual two-sample Bernoulli form rather than
the nested child-parent form used in the edge test. The standardized sibling
vector is
\begin{equation}
z_{u,j}^{\mathrm{sib}}
=
\frac{\theta_{l(u),j} - \theta_{r(u),j}}
{\sqrt{
\bar{\theta}_{u,j}(1-\bar{\theta}_{u,j})
\left(\frac{1}{n_{l(u)}} + \frac{1}{n_{r(u)}}\right)
\left(1+\frac{t_{u,l(u)}+t_{u,r(u)}}{2\bar t}\right)
}}.
\end{equation}
For categorical and continuous feature blocks, the implementation uses the same
two-sample sibling scale factor with the block covariance declared by the
feature-space contract: drop-last multinomial covariance for one-hot
categorical blocks and empirical-Gaussian covariance for continuous blocks.
The sibling test then asks how to combine this multi-feature contrast into one
summary score. The method reuses the spectral information already
computed during the child-parent stage rather than drawing a fresh random
projection from scratch. For a parent $u$ with two children $l(u)$ and
$r(u)$, let $k_{l(u)}^{\mathrm{edge}}$ and $k_{r(u)}^{\mathrm{edge}}$ denote
the positive edge-stage spectral dimensions of the two children. The sibling
dimension used in code is
\begin{equation}
k_u^{\mathrm{sib}}
=
\operatorname{round}
\!\left(
\sqrt{
k_{l(u)}^{\mathrm{edge}}
k_{r(u)}^{\mathrm{edge}}
}
\right).
\end{equation}
If exactly one child has a positive edge-stage spectral dimension, that value
is reused directly; if neither child has a positive edge-derived dimension, the
parent's own edge-stage spectral dimension is used. The sibling test is run only
for two-child sibling pairs with complete edge-stage spectral context. This
standardized vector answers the first sibling question: how large is the
left-right contrast after accounting for feature-specific sampling noise? The
next question is how to combine those coordinatewise contrasts into one summary
score.

\paragraph{Projected sibling statistic.}
The sibling test then combines this multi-feature contrast into a single
summary score while preserving the low-dimensional structure already learned at
the parent during the edge stage.

The projection basis starts with the parent node, not a fresh random draw. The
first available parent PCA directions, where PCA refers to principal-component
directions summarizing the main variation at that parent, are used up to
$k_u^{\mathrm{sib}}$. The requested sibling dimension is capped to the available
parent PCA rank, so the projection basis and eigenvalue vector have the same
dimension.

If $P_u \in \mathbb{R}^{k_u^{\mathrm{sib}} \times p}$ denotes this sibling
projection basis, the raw sibling statistic is
\begin{equation}
T_u^{\mathrm{raw}}
=
\|P_u z_u^{\mathrm{sib}}\|_2^2.
\end{equation}

\paragraph{Reference law.}
After projection, the method asks whether the observed squared contrast is
larger than would be expected from ordinary within-parent sampling noise. The
PCA eigenvalues select the parent-local subspace; they do not weight the null
law because the statistic is computed from inner products with orthonormal
projection rows. Thus, for a fixed selected subspace,
\[
T_u^{\mathrm{raw}} \mid H_0 \sim \chi^2_{k_u^{\mathrm{sib}}}.
\]

\paragraph{Worked sibling example.}
A small numerical example helps separate three ideas that can be conflated:
projection, per-component whitening, and the orthonormal projected-Wald
reference law. Suppose the standardized sibling difference vector is
\begin{equation}
z_u^{\mathrm{sib}} =
\begin{bmatrix}
2.0\\
-1.0\\
0.5\\
1.0
\end{bmatrix},
\qquad
P_u =
\begin{bmatrix}
1&0&0&0\\
0&1&0&0\\
0&0&1&0
\end{bmatrix}.
\end{equation}
Then the projected contrast is
\begin{equation}
P_u z_u^{\mathrm{sib}} =
\begin{bmatrix}
2.0\\
-1.0\\
0.5
\end{bmatrix}.
\end{equation}
Without any PCA variance correction, the raw quadratic form is
\begin{equation}
T_u^{\mathrm{raw}}
=
\|P_u z_u^{\mathrm{sib}}\|_2^2
=
2.0^2 + (-1.0)^2 + 0.5^2
=
5.25.
\end{equation}
Under the implemented orthonormal projected-Wald reference, this statistic is
compared directly to $\chi^2_3$. If one instead chose a different whitening
statistic, eigenvalues would enter before summation. For example, suppose the
first two projected directions have eigenvalues $\lambda_1 = 4$ and
$\lambda_2 = 1$, while the third coordinate is an added unit direction. In
matrix form,
\begin{equation}
\Lambda =
\begin{bmatrix}
4&0\\
0&1
\end{bmatrix},
\qquad
\Lambda^{-1} =
\begin{bmatrix}
1/4&0\\
0&1
\end{bmatrix}.
\end{equation}
Per-component whitening would then use
\begin{equation}
T_u^{\mathrm{white}}
=
\begin{bmatrix}
2.0 & -1.0
\end{bmatrix}
\Lambda^{-1}
\begin{bmatrix}
2.0\\
-1.0
\end{bmatrix}
+
0.5^2
=
\frac{2.0^2}{4}
+
\frac{(-1.0)^2}{1}
+
0.5^2
=
2.25.
\end{equation}
The first direction contributes less after whitening because a contrast of
size $2.0$ is less surprising along a direction whose null variance is already
$4$. The implemented sibling default does not use this whitened statistic. It
keeps the raw orthonormal projection, so the calibration is
\begin{equation}
T_u^{\mathrm{raw}} \sim \chi^2_3,
\qquad
p_u^{\mathrm{sib}}
=
\Pr\!\left(
\chi^2_3 \ge 5.25
\right).
\end{equation}
In general, the projected-Wald kernel records this unselected reference as
$a_u=1$ and $\nu_u=k_u^{\mathrm{sib}}$. \Cref{app:projected-wald-example}
expands this example in full matrix form.

\paragraph{Post-selection scale distortion.}
Raw sibling $p$-values are still not enough on their own because the sibling
pair being tested is not fixed in advance. The same data are used first to
construct the hierarchy and thereby expose pairs of child subtrees that already
look well separated, and then again to test whether those selected siblings are
significantly different. This data reuse creates \emph{post-selection
scale distortion}. Even under a null setting with no true cluster structure, the tree
construction step tends to present unusually separated sibling pairs because it
was built from the observed dissimilarities. The sibling statistic
is therefore evaluated conditionally on a favorable selection event rather than
under the unconditional null distribution for a pre-specified pair. The
practical consequence is anti-conservative inference: the observed sibling
difference vector $z_u^{\mathrm{sib}}$ can be larger than the naive reference
law expects, so raw $p$-values can be too small and false split decisions can
become too common. This motivation is consistent with the wider
post-selection-inference literature, where tests must account for the fact
that the reported hypothesis was selected using the data
\citep{fithian2014optimal,lee2016exact,gao2024selective}.
[RESULTS GAP: archived diagnostic files suggest that uncorrected sibling
testing can be anti-conservative, but the manuscript does not yet contain a
reproducible null simulation with sample size, random seeds, confidence
intervals, and the exact implemented correction. Do not report a numeric
Type~I error rate until that study is added.]

\paragraph{Context-weighted empirical-null inflation.}
The sibling calibration accounts for this extra optimism by estimating the
runtime empirical inflation of the projected-Wald statistic from the empirical
null visible inside the fitted hierarchy. The estimator has three ingredients:
a raw sibling statistic, a weight describing how null-like that sibling context
looks from the edge stage, and a feature-family-specific context used to
borrow information across similar projection dimensions and parent sizes.

For every two-child parent $q$, let $T_q^{\mathrm{raw}}$ be the raw sibling
statistic. Let $a_q$ and $\nu_q$ be the reference scale and degrees of freedom
returned by the projected-Wald kernel, with $a_q=1$ and $\nu_q=k_q$ for the
current orthonormal PCA basis. Let $\pi_q \in [0,1]$ be the sibling null weight
implied by the two child-parent edge tests. For each child edge, define
\(\omega_{q\to c}=1\) if the edge is not tested or is ancestor-blocked by the
Tree-BH path, and otherwise set \(\omega_{q\to c}\) to the finite Tree-BH
adjusted child-parent p-value. In the current implementation,
\begin{equation}
\pi_q =
\omega_{q\to l(q)}\,\omega_{q\to r(q)},
\end{equation}
so both child-parent edges must look empirically null-like for the sibling pair
to carry high calibration weight. This is a monotone empirical weight, not a
posterior null probability. Also,
\[
x_q = (\log s_q,\log n_q),
\]
where $s_q$ is the resolved sibling projection dimension and $n_q$ is the
parent leaf count. The empirical estimator uses the working post-selection scale family
\begin{equation}\label{eq:sibling-context-scale-model}
T_q^{\mathrm{raw}} \mid H_0
\approx
a_q\,c_{f_q}(x_q)\,\chi^2_{\nu_q}.
\end{equation}
Here $a_q$ is the analytic projected-quadratic reference scale. In the current
orthonormal projection implementation it is one, so $c_{f_q}(x_q)$ is the
empirical post-selection inflation for feature family $f_q$.

The estimator borrows calibration information most strongly from sibling
records with similar context values. For a focal sibling pair $u$, define
context weights
\begin{equation}
w_q(u)
=
\pi_q\,\mathbf 1\{f_q=f_u\}\,K_{f_u}(x_q,x_u),
\end{equation}
where the default kernel is Gaussian over the active context axes,
\begin{equation}
K_{f_u}(x_q,x_u)
=
\exp\!\left[
-\frac{1}{2}
\sum_{j\in A_{f_u}}
\left(
\frac{x_{qj}-x_{uj}}{h_j}
\right)^2
\right],
\end{equation}
where $A_{f_u}$ is the set of active context axes with nonzero empirical
bandwidth. The current implementation always uses \(\log s_q\) when it has
positive bandwidth and additionally uses \(\log n_q\) for non-Bernoulli
feature families when that axis has positive bandwidth. The local
empirical-null inflation estimate is the weighted
chi-square scale maximum-likelihood estimator after accounting for the analytic
projected-Wald reference scale:
\begin{equation}\label{eq:sibling-context-scale-estimator}
\hat{c}_{f_u}(x_u)
=
\max\left\{
1,\,
\frac{\sum_q w_q(u)T_q^{\mathrm{raw}}}
{\sum_q w_q(u)a_q\nu_q}
\right\}.
\end{equation}
Let $g(q)$ identify the stopping event that owns calibration record $q$.
Nested blocked descendants inherit the group of the nearest tested,
non-significant ancestor edge path. Support counts, effective sample size,
maximum weight share, and leave-one-group sensitivity are evaluated over these
dependency groups, so descendants of one stopping event cannot create
additional effective support. These quantities remain support diagnostics;
grouping known shared events does not establish independence between groups or
validate a selected-tail reference law.

\paragraph{Calibration-support contract.}
For a focal sibling pair $u$, the empirical-null interpretation of
\(\hat c_{f_u}(x_u)\) requires positive calibration mass from records that are
tested null-like or lie below a stopped edge path.
Let $f_q$ denote the feature family of sibling record $q$, let
$E_{l(q)}^{\mathrm{edge}}$ and $E_{r(q)}^{\mathrm{edge}}$ indicate whether the
two child-parent edges were tested, let $S_{l(q)}^{\mathrm{edge}}$ and
$S_{r(q)}^{\mathrm{edge}}$ be their rejection indicators, and let
$B_q^{\mathrm{edge}}$ indicate that the sibling record is edge-blocked by the
ancestor testing path. Define the local support set
\[
\mathcal C_u
=
\left\{
q:
f_q=f_u,\ \nu_q>0,\ w_q(u)>0
\right\}.
\]
The strict empirical-null support set is
\[
\mathcal C_{0,u}
=
\left\{
q\in\mathcal C_u:
E_{l(q)}^{\mathrm{edge}}=1,\
E_{r(q)}^{\mathrm{edge}}=1,\
S_{l(q)}^{\mathrm{edge}}=0,\
S_{r(q)}^{\mathrm{edge}}=0
\right\},
\]
so untested descendants are not called strict-null observations. The stopped
support set is
\[
\mathcal C_{\mathrm{stop},u}
=
\left\{
q\in\mathcal C_u:
B_q^{\mathrm{edge}}=1
\right\}.
\]
If $\mathcal C_{0,u}$ has positive calibration mass, the local estimate has
strict internal empirical-null support. If $\mathcal C_{0,u}$ is empty but
$\mathcal C_{\mathrm{stop},u}$ has positive mass, the estimate has only
stopped-path support. The current estimator retains nested blocked descendants
as node records but assigns them to their owning stopping-event dependency
group. If both support sets are empty but $\mathcal C_u$ has positive
mass, the data provide selected non-null context but not empirical-null
calibration. If $\mathcal C_u$ has no positive mass, the local empirical-null
calibration contract is unsatisfied. The implementation reports
leave-one-group sensitivity, but its configured thresholds are not validated
method constants.

The production rule is intentionally strict. The current implementation raises
a calibration-data error when the available records are only selected non-null
contexts. The current diagnostics do not justify a separately named external
calibration model for this case; production therefore fails closed. Any future
external selection-conditioned model would first require validation under the
full hierarchy-selection event, rather than being used as a fallback. The method
must not silently substitute a neutral inflation factor.
The present evidence establishes this as a necessary support contract, not as a
sufficient finite-sample guarantee. In particular, internal support does not
resolve the selected-hierarchy reference law.
The denominator is required to have positive effective calibration mass from
tested-null or stopped-path records. If no such positive-degree sibling
records carry positive null weight, the calibration contract is not satisfied
and the pipeline closes focal sibling gates; the benchmark reports an unsupported outcome.

\paragraph{Reference-law contract and sibling false-discovery-rate correction.}
For same-data selected hierarchies, the hierarchy, focal pair, calibration
records, stopping events, and weights are all data-derived and may share
observations. Consequently,
\begin{equation}
\Pr\!\left(
\chi^2_{\nu_u}
\ge
\frac{T_u^{\mathrm{raw}}}{a_u\hat c_{f_u}(x_u)}
\right)
\end{equation}
is the restored empirical plug-in p-value, whose selected-tail validity remains
unproven. With internal support, the production path reports
\texttt{internal\_admissible} and supplies this p-value to traversal-aligned
sibling BH. Here admissibility means availability under the support policy,
not a selective-inference guarantee. Dependency-group support diagnostics
and the unresolved reference-law metadata remain available in the decision API.
The zero-dimensional case remains the deterministic decision $p=1$.

A separate exact mode is available only under a stronger contract. Suppose a
focal statistic and $m$ calibration statistics satisfy
\[
\frac{T_u}{ac}\sim\chi^2_{\nu_u},
\qquad
\frac{T_q}{ac}\sim\chi^2_{\nu_q},
\]
mutually independently, with common positive reference scale $a$, common
unknown inflation $c$, and fixed unit calibration weights. Let
$D=\sum_{q=1}^m\nu_q$ and
\[
\widetilde c=\frac{\sum_{q=1}^m T_q}{aD}.
\]
Pairwise-disjoint calibration observation sets and a focal observation set
disjoint from their union are required provenance for this mode. Then
\[
F_u=\frac{T_u}{a\nu_u\widetilde c}\sim F_{\nu_u,D}.
\]
The exact-mode p-value is the upper $F_{\nu_u,D}$ tail, floored by the
unadjusted upper $\chi^2_{\nu_u}$ tail so estimated deflation cannot make the
one-sided calibrated p-value smaller. Unequal deterministic weights are not
accepted by this exact mode; neither Kish effective sample size nor a
Satterthwaite approximation is labeled exact.

The sibling multiplicity correction is ordinary
Benjamini--Hochberg \citep{benjamini1995controlling} at
$\alpha_{\mathrm{sib}} = 0.01$, applied only across focal sibling pairs. Let
$S^{\mathrm{sib}}(u) \in \{0,1\}$ denote the resulting indicator that parent
$u$ supports a significant sibling split.
